% =============================================================================
%  Chapter 4 — Gauss-Point to Nodal Projection Methods
% =============================================================================

\chapter{Gauss-Point to Nodal Projection}
\label{ch:projection}

In displacement-based finite elements, the primary unknowns (displacements)
are nodal.  Derived quantities\,---\,strains, stresses, internal
variables\,---\,are computed at \emph{Gauss (integration) points} inside each
element.  For post-processing and visualisation, these quantities must be
\emph{projected} or \emph{averaged} to the mesh nodes so that they can be
rendered as continuous fields in VTK/ParaView.

This chapter compares the three projection methods implemented in
\texttt{fall\_n} and explains the current adaptive default: lumped $L^2$ when
the local nodal lumping is strictly positive, and volume-weighted averaging
otherwise.


% ─────────────────────────────────────────────────────────────────────────
\section{Lumped \texorpdfstring{$L^2$}{L2} Projection}
\label{sec:l2_projection}

\subsection{Consistent \texorpdfstring{$L^2$}{L2} Projection}

The mathematically rigorous approach minimises the $L^2$ norm of the error
between the discrete Gauss-point field and a continuous nodal field.  Let
$\sigma^h$ be the Gauss-point field and $\tilde{\sigma}^h = \sum_I
\tilde{\sigma}_I N_I$ the nodal projection.  The consistent $L^2$ projection
solves:
\begin{equation}
  \mathbf{M}\,\tilde{\bsigma} = \mathbf{f},
  \label{eq:consistent_L2}
\end{equation}
where $\mathbf{M}$ is the global mass matrix ($M_{IJ} = \int_\Omega N_I N_J
\,\dd\Omega$) and $\mathbf{f}$ is the right-hand side ($f_I = \int_\Omega N_I
\sigma^h \,\dd\Omega$).

This requires solving a global linear system, which is expensive and may
introduce oscillations at element boundaries (Gibbs-like effects) for
discontinuous fields.

\subsection{Row-Sum Lumping}
\label{sec:lumping}

The standard shortcut is to \emph{lump} the mass matrix by replacing it with
its row sums on the diagonal:
\begin{equation}
  \tilde{M}_{II} = \sum_J M_{IJ}
  = \int_\Omega N_I \sum_J N_J \,\dd\Omega
  = \int_\Omega N_I\,\dd\Omega,
  \label{eq:lumped_mass}
\end{equation}
using the partition of unity $\sum_J N_J = 1$.  The lumped projection then
becomes a local (node-by-node) average:
\begin{equation}
  \tilde{\sigma}_I
  = \frac{\displaystyle\int_\Omega N_I\, \sigma^h \,\dd\Omega}
         {\displaystyle\int_\Omega N_I \,\dd\Omega}
  \approx
  \frac{\displaystyle\sum_{e \in \mathcal{E}_I}
        \sum_{g=1}^{n_g} N_I(\bxi_g)\, w_g\, |J_g|\, \sigma_g}
       {\displaystyle\sum_{e \in \mathcal{E}_I}
        \sum_{g=1}^{n_g} N_I(\bxi_g)\, w_g\, |J_g|}.
  \label{eq:lumped_L2_final}
\end{equation}

The denominator is the ``nodal volume'' contributed by the quadrature rule.
For it to be meaningful, we need
\begin{equation}
  D_I = \sum_{e \in \mathcal{E}_I}
        \sum_g N_I(\bxi_g)\, w_g\, |J_g| > 0
        \qquad \forall\, I.
  \label{eq:denominator_positive}
\end{equation}

As demonstrated in \cref{ch:negative_weights} and refined in
\cref{ch:adaptive_projection}, this condition can fail in two distinct ways:
\begin{itemize}[nosep]
  \item \emph{catastrophically}, when negative quadrature weights amplify the
        signed support of higher-order simplex basis functions;
  \item \emph{structurally}, when the exact row-sum lumping
        $\int_\Omega N_I\,\dd\Omega$ is itself non-positive, as happens for the
        vertex functions of TET10 even under all-positive quadrature.
\end{itemize}
In the second case the formula is still well-defined, but it is no longer a
nodal average.  It becomes a signed extrapolation operator, which is precisely
what produces the spurious oscillations seen in TET10 contour plots.

\paragraph{Implementation.}  The lumped $L^2$ projection is implemented in
\texttt{l2\_project\_to\_nodes()} in \texttt{VTKModelExporter.hh}.  The legacy
\texttt{std::abs} guard on the denominator remains in the explicit
\texttt{LumpedL2} path as a defensive measure against future negative-weight
rules:
\begin{lstlisting}
double denom = std::abs(node_lump[I]);
if (denom < 1e-14) denom = 1e-14;
node_val[I] /= denom;
\end{lstlisting}
but this is no longer relied on as the production fix for quadratic simplex
elements.  The production path is now the adaptive resolver documented in
\cref{ch:adaptive_projection}.


% ─────────────────────────────────────────────────────────────────────────
\section{Volume-Weighted Averaging (SPR-Type)}
\label{sec:spr_averaging}

\subsection{Motivation}

The Superconvergent Patch Recovery (SPR) method of Zienkiewicz and
Zhu~\cite{ZienkiewiczZhu1992a,ZienkiewiczZhu1992b} fits a polynomial to the
Gauss-point values in a \emph{patch} of elements surrounding each node and
evaluates the polynomial at the node.  In its full form, SPR requires solving
local least-squares problems.

For production post-processing where visual smoothness (rather than
superconvergent error estimation) is the primary goal, a simpler and even more
robust approach suffices: \textbf{volume-weighted element averaging}.

\subsection{Algorithm}
\label{sec:spr_algorithm}

For each node $I$, compute the nodal field as a weighted average of the
element centroid values:
\begin{equation}
  \boxed{
  \tilde{\sigma}_I
  = \frac{\displaystyle\sum_{e \in \mathcal{E}_I} V_e\, \bar{\sigma}_e}
         {\displaystyle\sum_{e \in \mathcal{E}_I} V_e},
  }
  \label{eq:spr_average}
\end{equation}
where:
\begin{itemize}[nosep]
  \item $\mathcal{E}_I$ is the set of elements sharing node $I$;
  \item $V_e$ is the volume of element $e$ (computed via quadrature:
        $V_e = \sum_g w_g |J_g|$);
  \item $\bar{\sigma}_e$ is the element-average value, computed as the
        weighted mean of Gauss-point values:
        \begin{equation}
          \bar{\sigma}_e
          = \frac{\sum_g w_g\, |J_g|\, \sigma_g}{\sum_g w_g\, |J_g|}.
          \label{eq:element_average}
        \end{equation}
\end{itemize}

\subsection{Robustness Guarantee}

\begin{proposition}[Unconditional positivity]
All weights in the averaging formula~\eqref{eq:spr_average} are element
volumes $V_e > 0$, which are strictly positive for any non-degenerate element
regardless of the quadrature rule used for stiffness integration.  Even if the
quadrature rule contains negative weights, the element volume $V_e = \sum_g
w_g |J_g|$ remains positive (it is the exact volume when the quadrature degree
is sufficient to integrate polynomials of the Jacobian's degree).

Therefore, SPR volume-weighted averaging is \textbf{unconditionally robust}:
the denominator is always positive, and no sign flip or division by zero can
occur at any node, regardless of element type or quadrature rule.
\end{proposition}

\subsection{Comparison with Lumped \texorpdfstring{$L^2$}{L2}}

\begin{center}
\renewcommand{\arraystretch}{1.2}
\begin{tabular}{p{4cm}cc}
  \toprule
  Property & Lumped $L^2$ & Volume-weighted avg. \\
  \midrule
  Denominator sign    & Can be $\le 0$ & Always $> 0$ \\
  Sensitivity to $w_g < 0$ & High (catastrophic) & None \\
  Accuracy order      & Optimal ($O(h^{k+1})$) & $O(h^k)$ \\
  Superconvergent?    & No  & No \\
  Computational cost  & $O(n_e \cdot n_g)$ & $O(n_e \cdot n_g)$ \\
  Requires global solve? & No & No \\
  \bottomrule
\end{tabular}
\end{center}

For visualisation purposes, the $O(h^k)$ vs.\ $O(h^{k+1})$ accuracy
difference is negligible.  The robustness advantage of volume-weighted
averaging is decisive.


% ─────────────────────────────────────────────────────────────────────────
\section{Implementation: \texttt{spr\_average\_to\_nodes()}}
\label{sec:spr_impl}

The volume-weighted averaging is implemented in two methods in
\texttt{VTKModelExporter.hh}:

\begin{description}
  \item[\texttt{spr\_average\_to\_nodes()}]
    Computes a single scalar field $\sigma$:
    \begin{enumerate}[nosep]
      \item For each element $e$, compute $V_e$ and $\bar{\sigma}_e$ from
            Gauss-point data.
      \item For each node $I$ of element $e$, accumulate:
            \begin{align}
              \text{numerator}[I]   &\mathrel{+}= V_e \cdot \bar{\sigma}_e, \\
              \text{denominator}[I] &\mathrel{+}= V_e.
            \end{align}
      \item After all elements: $\tilde{\sigma}_I = \text{numerator}[I] \,/\,
            \text{denominator}[I]$.
    \end{enumerate}

  \item[\texttt{compute\_material\_fields\_spr()}]
    Applies \texttt{spr\_average\_to\_nodes()} to all material fields: strain
    tensor (6 components), stress tensor (6 components), plastic strain tensor
    (6 components), and equivalent plastic strain (1 scalar).

  \item[\texttt{compute\_material\_fields()}]
    Resolves the projection strategy automatically.  It uses lumped $L^2$
    only if every element in the model has strictly positive local lumped nodal
    weights; otherwise it dispatches to volume-weighted averaging.
\end{description}


% ─────────────────────────────────────────────────────────────────────────
\section{Verification}
\label{sec:projection_verification}

The adaptive projection policy was verified on the cantilever benchmark
($10 \times 0.4 \times 0.8$ box, $E = 200$, $\nu = 0.3$):
\begin{itemize}[nosep]
  \item \textbf{HEX27 mesh} (20 elements, 315 nodes): the local lumped nodal
        weights are strictly positive, so the adaptive policy resolves to
        lumped $L^2$.  The projected fields remain smooth.
  \item \textbf{TET10 mesh} ($\sim 1329$ elements, 2840 nodes):
    \begin{itemize}[nosep]
      \item Lumped $L^2$ with Keast-5: \textbf{checkerboard at vertices}.
      \item Lumped $L^2$ with HMS-4 or Stroud conical product: still not a
            true nodal average, because the TET10 vertex lump weights remain
            negative.
      \item Adaptive mode therefore resolves to SPR volume-weighted
            averaging, which yields smooth, correct fields.
    \end{itemize}
  \item The main showcase (\texttt{main.cpp}) now uses
        \texttt{compute\_material\_fields()} with adaptive dispatch for all
        14~cases (7~formulations $\times$ 2~element types).
\end{itemize}
